% Clase 1 --- Ley de Coulomb en forma vectorial (§3.3).
% Dos paneles con la MISMA geometría (Q2 arriba, Q1 abajo) para que se vea que
% lo único que cambia es el signo del producto Q1 Q2, y con él el sentido de F12.
% El versor r21 va SIEMPRE de 2 hacia 1; F12 es la fuerza que 2 le hace a 1.
% F12 se dibuja en una vertical desplazada para no superponerse con r21.
\begin{center}
\begin{tikzpicture}[scale=1]

  % =============== panel (a): mismo signo -> repulsión ===============
  \begin{scope}
    \node[figpos] (a2) at (0,2.6) {};
    \node[figpos] (a1) at (0,0)   {};
    \node[figlbl, left=4pt] at (a2) {$Q_2 > 0$};
    \node[figlbl, left=4pt] at (a1) {$Q_1 > 0$};
    % recta que une las cargas
    \draw[figaux] (0,3.15) -- (0,-1.7);
    % vector r21 (de 2 a 1), sobre la recta que une las cargas
    \draw[figvecaux] (a2) -- (a1);
    \node[figlblaux, left=2pt] at (0,1.45) {$\vec r_{21}$};
    % versor r21
    \draw[figvec, figgray] (1.9,2.6) -- (1.9,1.75);
    \node[figlblaux, right=2pt] at (1.9,2.2) {$\hat r_{21}$};
    % fuerza sobre 1: mismo sentido que el versor
    \draw[figvecres] (0.55,0) -- (0.55,-1.4);
    \node[figlbl, figred, right=3pt] at (0.55,-0.75) {$\vec F_{12}$};
    \node[figlbl, align=center] at (0.4,-2.35)
      {(a) $Q_1Q_2 > 0$\\[1pt]
       $\vec F_{12}$ va como $\hat r_{21}$: \textbf{repulsión}};
  \end{scope}

  % =============== panel (b): signos opuestos -> atracción ===============
  \begin{scope}[xshift=7.2cm]
    \node[figpos] (b2) at (0,2.6) {};
    \node[figneg] (b1) at (0,0)   {};
    \node[figlbl, left=4pt] at (b2) {$Q_2 > 0$};
    \node[figlbl, left=4pt] at (b1) {$Q_1 < 0$};
    \draw[figaux] (0,3.15) -- (0,-1.7);
    \draw[figvecaux] (b2) -- (b1);
    \node[figlblaux, left=2pt] at (0,1.45) {$\vec r_{21}$};
    \draw[figvec, figgray] (1.9,2.6) -- (1.9,1.75);
    \node[figlblaux, right=2pt] at (1.9,2.2) {$\hat r_{21}$};
    % fuerza sobre 1: sentido opuesto al versor
    \draw[figvecres] (0.55,0) -- (0.55,1.4);
    \node[figlbl, figred, right=3pt] at (0.55,0.75) {$\vec F_{12}$};
    \node[figlbl, align=center] at (0.4,-2.35)
      {(b) $Q_1Q_2 < 0$\\[1pt]
       $\vec F_{12}$ va como $-\hat r_{21}$: \textbf{atracción}};
  \end{scope}

\end{tikzpicture}\\[0.5em]
{\small\itshape Una sola expresión contiene los dos casos: el versor
$\hat r_{21}$ fija la dirección (el eje que une las cargas) y el signo del
producto $Q_1Q_2$ fija el sentido. La regla ``iguales se repelen, opuestas se
atraen'' ya está dentro de la fórmula, no hay que agregarla aparte.}
\end{center}
